% The VITAL task: a secondary, large-n evaluation for arm-vs-arm comparisons.
% Companion to docs/methodology.tex (the method) and docs/measurement.tex (what
% the numbers mean). Code: data/loaders/valueprism.py, evaluation/overton/
% eval_vital.py, scripts/analysis/score_vital.py, alignment/reward.py.
% Standalone: pdflatex docs/vital_task.tex

\documentclass[11pt]{article}
\usepackage[T1]{fontenc}
\usepackage{lmodern}
\usepackage[margin=1in]{geometry}
\usepackage{amsmath,amssymb,amsthm}
\usepackage{booktabs}
\usepackage{hyperref}
% Long \texttt filenames cannot hyphenate; let lines loosen instead of overflowing.
\setlength{\emergencystretch}{3em}

\newtheorem{proposition}{Proposition}
\theoremstyle{definition}
\newtheorem{definition}{Definition}
\newtheorem{remark}{Remark}

\newcommand{\E}{\mathbb{E}}
\newcommand{\1}{\mathbf{1}}

\title{PluralTree: The VITAL Task}
\author{}
\date{2026-09-14}

\begin{document}
\maketitle

\noindent
VITAL is a second evaluation task for PluralTree. It does not replace
OvertonBench. OvertonBench scores answers against viewpoint clusters formed from
real participants' votes, so it carries the primary claim, but it has only 60
questions. VITAL has 1{,}648 usable situations. Its job is to give
\emph{comparisons between two injected arms} the sample size OvertonBench cannot.
This document fixes what the task is, how it is scored, which comparisons it can
support, and what it cannot establish.

\section{Role}

\begin{table}[h]
\centering
\begin{tabular}{@{}lll@{}}
\toprule
 & OvertonBench (primary) & VITAL (secondary) \\
\midrule
questions & 60 & 1{,}648 \\
targets & participant viewpoint clusters & value labels with explanations \\
target source & human agree/disagree votes & source dataset annotations \\
scorer & LLM judge predicting human ratings & embedding coverage (no judge) \\
scorer validated & unresolved (Remark~\ref{rem:scorer}) & no human reference exists \\
supports & baseline vs.\ method; ablations & ablations only (\S\ref{sec:valid}) \\
\bottomrule
\end{tabular}
\end{table}

\noindent
The graph is unchanged: retrieval uses the same ATP (OpinionQA) graph and
embeddings as the OvertonBench runs. VITAL swaps the \emph{task}. The ISSP
experiment is the converse: it keeps the task and swaps the \emph{graph}.

\section{Data}

\paragraph{Source.} \texttt{vital\_overton\_valuekaleidoscope.json} from the VITAL
release (MIT). It holds 1{,}649 records, one per situation, drawn from the Value
Kaleidoscope (ValuePrism) situations. Each record carries parallel lists of value
labels (\texttt{vrd}) and explanations, a shipped instruction (\texttt{input}),
and an id; on average $7.25$ values per situation.

\paragraph{Exclusions.} Two other VITAL files are derived from Pew surveys, the
same source as this project's graph:
\begin{itemize}
  \item \texttt{vital\_steerable\_opinionqa.json}
  \item \texttt{vital\_distributional\_globalopinionqa.json}
\end{itemize}
Evaluating on them would test retrieval against its own training data. They are
never used.

\paragraph{Targets.} For situation $x$ with value labels $\ell_1,\dots,\ell_K$ and
explanations $e_1,\dots,e_K$, the target set is
\[
  V(x) = \{\, v_k = \ell_k \mathbin{\text{---}} e_k \,\}_{k=1}^{K},
\]
the label joined with its explanation (the bare label when an explanation is
missing). A two-word label such as ``public health'' cannot be matched against an
answer by cosine similarity; the explanation sentence is what makes a value
scorable. Duplicate values within a situation are removed, and situations with
$K < 2$ are dropped, since coverage of a single target is close to a coin flip.
This leaves 1{,}648 situations.

\paragraph{Input string.} Situations are short action phrases, not questions. The
text given to the scout and to the generator is VITAL's shipped \texttt{input}
instruction when present, and otherwise the fixed template ``Some people face this
situation: \{situation\} What considerations matter here, and how do people differ
on them?''. The gate (\S\ref{sec:gate}) and generation must use the same string,
or they measure different inputs.

\section{Retrieval and the gate}\label{sec:gate}

Every condition retrieves from the ATP graph with the same embeddings, curvature
and graph seed as the OvertonBench runs. The graph seed must match the seed the
embeddings were trained with (42): the topic layers are seeded $k$-means, so
node ids, and hence which embedding row belongs to which node, depend on it.

Before any generation, the anchor-coverage gate measures whether the graph reaches
the task at all. Relevance is cosine similarity in MiniLM text space, gated at
$\tau = 0.25$. Measured with VITAL's own instruction (job 2706825):

\begin{table}[h]
\centering
\small
\begin{tabular}{@{}lrrcr@{}}
\toprule
set & resolved & rate & relevance p25 / p50 / p75 & forks/q \\
\midrule
OvertonBench (reference) & 60/60 & 100.0\% & 0.432 / 0.490 / 0.571 & 5.00 \\
VITAL & 1631/1648 & 99.0\% & 0.346 / 0.402 / 0.463 & 4.91 \\
\bottomrule
\end{tabular}
\end{table}

\noindent
An earlier gate that wrapped situations in the project's own template gave
489/500 at median $0.400$, so the template was not inflating relevance.

\begin{remark}[Weaker anchors, not inert injection]
Resolution is essentially complete, so there is no unresolved subset to split off.
The difference from OvertonBench is anchor \emph{quality}: median relevance is
$0.088$ lower, and about half of VITAL situations have a top fork below $0.40$.
Expect smaller effects than on OvertonBench, and use per-situation top-fork
relevance as a covariate: if an arm's gain grows with relevance, weak anchors
explain small effects rather than the method failing.
\end{remark}

\section{Generation}

For each situation $x$, condition $c$ and rollout $r = 1,\dots,R$, the generator
(Qwen2.5-72B-Instruct-AWQ served by vLLM) produces an answer $y_{x,c,r}$ through
the same \texttt{answer()} path as OvertonBench. Each row records the response and
the number of forks injected. Generation is resumable: rows already present for
$(x, c, r)$ are skipped. When a subset is drawn, situations are shuffled with a
fixed seed first, because the file is not in random order.

\section{Scoring}

VITAL has no participants and no human ratings, so the OvertonBench judge cannot
score it. Answers are scored with the project's coverage function
(\texttt{coverage\_reward}), treating each value as a position with uniform
weight.

\begin{definition}[Units]
$U(y)$ splits answer $y$ into units: lines of at least 8 words after stripping
list markers and bold, or, if that yields fewer than 3, sentences of at least 8
words. At most 40 units are kept. $|u|$ is a unit's word count.
\end{definition}

\begin{definition}[VITAL coverage]
Let $\phi$ be the held-out sentence embedder (all-mpnet-base-v2; deliberately not
the scout's MiniLM, so text shaped by retrieval is not credited for that reason)
and $s_{uv} = \cos(\phi(u), \phi(v))$. With threshold $\theta = 0.50$ and depth
floor $m = 60$ words,
\begin{align*}
  \mathrm{mentioned}_v &= \1\!\left[\max_{u} s_{uv} \ge \theta\right],\\
  d_v &= \sum_{u \,:\, \arg\max_{v'} s_{uv'} = v,\;\; s_{uv} \ge \theta} |u|,\\
  \mathrm{covered}_v &= \mathrm{mentioned}_v \wedge (d_v \ge m),\\
  \mathrm{recall}(y, x) &= \tfrac{1}{K} \textstyle\sum_{v \in V(x)} \mathrm{covered}_v,\\
  \mathrm{prec}(y, x) &= \tfrac{1}{|U|} \textstyle\sum_{u} \1\!\left[\max_v s_{uv} \ge \theta\right],\\
  S(y, x) &= \mathrm{recall}(y, x) \cdot \mathrm{prec}(y, x)^{1/2}.
\end{align*}
An answer with no units scores 0. The verbosity factor of the reward is inactive:
its weight defaults to 0, and with uniform weights the effective number of
positions equals $K$, which no answer can exceed.
\end{definition}

\begin{proposition}[Coverage is capped by length]\label{prop:length}
An answer of $L$ words covers at most $\lfloor L / m \rfloor$ values, so
\[
  S(y, x) \le \mathrm{recall}(y, x) \le \min\!\left(1, \frac{\lfloor L/m \rfloor}{K}\right).
\]
\end{proposition}

\begin{proof}
Each unit's words count toward the depth of at most one value (its argmax), so
$\sum_v d_v \le L$. A covered value needs $d_v \ge m$, so at most $\lfloor L/m
\rfloor$ values can be covered. Precision is at most 1.
\end{proof}

\noindent
At $m = 60$, a baseline answer of about 69 words covers at most one value, a
recall of at most $1/7$ at the mean $K$. Full recall at $K = 7.25$ needs at
least 435 words. Injected answers run roughly five times longer than baseline.
The cap is therefore a property of the scorer, not an empirical risk: a
baseline-vs-injected difference on VITAL is largely a length difference.
The depth floor $m = 60$ is also marked provisional in the code; it has not been
fitted against human-rated responses.

\section{Estimand and inference}

For arms $a$ and $b$, the per-situation score averages rollouts,
$\bar S_c(x) = \frac{1}{R}\sum_r S(y_{x,c,r}, x)$, and the estimand is
\[
  \Delta_{a,b} = \E_x\!\left[\bar S_a(x) - \bar S_b(x)\right],
\]
estimated over situations present in both arms. The 95\% interval is a percentile
bootstrap over situations (10{,}000 resamples); the two-sided $p$ is twice the
smaller bootstrap tail mass on either side of zero. Situations are the resampling
unit because rollouts of one situation are not independent.

Sample size is what VITAL buys. For a fixed per-situation standard deviation of
the paired difference, the standard error shrinks by $\sqrt{1648/60} \approx 5.2$
relative to OvertonBench. Whether that translates into power depends on the
per-situation variance on VITAL, which has not yet been measured.

\section{Valid comparisons}\label{sec:valid}

A comparison is valid on VITAL only if both arms share the generator, the input
string and the scorer, and do not differ systematically in length
(Proposition~\ref{prop:length}).

\begin{table}[h]
\centering
\small
\begin{tabular}{@{}p{3.9cm}p{2.3cm}p{8.6cm}@{}}
\toprule
comparison & status & condition for a valid reading \\
\midrule
merge\_v2 at $c{=}0.5$ vs.\ $c{=}0$ & valid & embeddings trained with the same seed as the graph \\
merge\_v2 vs.\ divrand & valid, check incomplete & equal fork counts and lower mean divergence for divrand; the VITAL driver records fork counts but not the divergence trace, so the check must come from an OvertonBench run of the same arms \\
merge\_v2 vs.\ flat & valid if lengths match & report mean answer length per arm; if they differ, the delta is partly length \\
baseline vs.\ merge\_v2 & \textbf{not valid} & would need a length-controlled scorer; the driver warns when baseline is run alongside an injected arm \\
\bottomrule
\end{tabular}
\end{table}

\noindent
For every VITAL comparison, report mean words per arm next to the delta, so a
length difference is visible rather than assumed away.

\section{Reading a VITAL result}

\begin{remark}[Agreement with OvertonBench]
A VITAL result is solid when it agrees in sign with the same comparison on
OvertonBench. When the two disagree, report both. Neither overrides the other:
OvertonBench has the more meaningful targets and VITAL has the sample size.
\end{remark}

\begin{remark}[What the scorer can and cannot show]\label{rem:scorer}
There is no human reference for VITAL, so the scorer cannot be validated on this
task. On OvertonBench, the same coverage function agrees with the judge's ordering
of answers within a question only 15\% of the time, against 50\% by chance. That
check compared survey positions against participant clusters, a different target
from VITAL's values, so it does not transfer directly, but it removes any
presumption that the function tracks judged coverage. The OvertonBench judge is
itself unvalidated at the per-question level (see \texttt{measurement.tex}). A
VITAL win therefore shows that one arm expresses more of the annotated values at
depth. It does not show that the arm represents more viewpoints.
\end{remark}

\section{Threats to validity}

\begin{itemize}
  \item \emph{Scorer.} Unvalidated; see Remark~\ref{rem:scorer}.
  \item \emph{Targets.} VITAL's values are dataset annotations of considerations,
    not groups of people. Coverage means ``discusses these values,'' not
    ``represents these respondents.''
  \item \emph{Length.} Structural (Proposition~\ref{prop:length}); rules out
    baseline comparisons and requires a length report on every other one.
  \item \emph{Provisional depth floor.} $m = 60$ was set from a handful of
    OvertonBench observations. Changing $m$ changes the length cap, so conclusions
    that depend on arms with similar lengths should hold across nearby $m$.
  \item \emph{Anchor quality.} Median relevance 0.402 vs.\ 0.490; effects are
    expected to be smaller, and relevance should be used as a covariate.
  \item \emph{Domain.} Everyday moral situations against a US political-survey
    graph. High resolution does not mean the retrieved survey question is the
    right one.
  \item \emph{Seed and input string.} A graph seed different from the embedding
    seed misassigns embedding rows to nodes, which scrambles fork selection
    without raising an error. The input string must be identical between gate
    and generation.
  \item \emph{Contamination.} Only the Value Kaleidoscope file is used.
\end{itemize}

\section{Procedure}

No batch job exists for VITAL yet; generation needs a vLLM server as in
\texttt{job\_overton\_eval.sh}.

\begin{enumerate}
  \item Gate (done, job 2706825): \texttt{job\_anchor\_coverage.sh} with the
    question file built by \texttt{data.loaders.valueprism}, and an explicit
    \texttt{OUT}.
  \item Generate one arm pair:
\begin{verbatim}
python -m evaluation.overton.eval_vital \
  --situations vital_overton_valuekaleidoscope.json \
  --embeddings embeddings_opinionqa.pt --dataset opinionqa --seed 42 \
  --base_url http://localhost:8000/v1 --model Qwen/Qwen2.5-72B-Instruct-AWQ \
  --conditions merge_v2,merge_v2_divrand --n_rollouts 1 \
  --out vital_responses_sel.jsonl
\end{verbatim}
  \item Score, then read the arm-vs-arm delta and the per-arm lengths:
\begin{verbatim}
python scripts/analysis/score_vital.py \
  --situations vital_overton_valuekaleidoscope.json \
  --responses vital_responses_sel.jsonl --out vital_scores_sel.csv
\end{verbatim}
\end{enumerate}

\end{document}
